% Options for packages loaded elsewhere
\PassOptionsToPackage{unicode}{hyperref}
\PassOptionsToPackage{hyphens}{url}
\PassOptionsToPackage{dvipsnames,svgnames,x11names}{xcolor}
\documentclass[
  11pt,
]{article}
\usepackage{xcolor}
\usepackage[margin=2.5cm]{geometry}
\usepackage{amsmath,amssymb}
\setcounter{secnumdepth}{-\maxdimen} % remove section numbering
\usepackage{iftex}
\ifPDFTeX
  \usepackage[T1]{fontenc}
  \usepackage[utf8]{inputenc}
  \usepackage{textcomp} % provide euro and other symbols
\else % if luatex or xetex
  \usepackage{unicode-math} % this also loads fontspec
  \defaultfontfeatures{Scale=MatchLowercase}
  \defaultfontfeatures[\rmfamily]{Ligatures=TeX,Scale=1}
\fi
\usepackage{lmodern}
\ifPDFTeX\else
  % xetex/luatex font selection
\fi
% Use upquote if available, for straight quotes in verbatim environments
\IfFileExists{upquote.sty}{\usepackage{upquote}}{}
\IfFileExists{microtype.sty}{% use microtype if available
  \usepackage[]{microtype}
  \UseMicrotypeSet[protrusion]{basicmath} % disable protrusion for tt fonts
}{}
\makeatletter
\@ifundefined{KOMAClassName}{% if non-KOMA class
  \IfFileExists{parskip.sty}{%
    \usepackage{parskip}
  }{% else
    \setlength{\parindent}{0pt}
    \setlength{\parskip}{6pt plus 2pt minus 1pt}}
}{% if KOMA class
  \KOMAoptions{parskip=half}}
\makeatother
\usepackage{longtable,booktabs,array}
\newcounter{none} % for unnumbered tables
\usepackage{calc} % for calculating minipage widths
% Correct order of tables after \paragraph or \subparagraph
\usepackage{etoolbox}
\makeatletter
\patchcmd\longtable{\par}{\if@noskipsec\mbox{}\fi\par}{}{}
\makeatother
% Allow footnotes in longtable head/foot
\IfFileExists{footnotehyper.sty}{\usepackage{footnotehyper}}{\usepackage{footnote}}
\makesavenoteenv{longtable}
\usepackage{graphicx}
\makeatletter
\newsavebox\pandoc@box
\newcommand*\pandocbounded[1]{% scales image to fit in text height/width
  \sbox\pandoc@box{#1}%
  \Gscale@div\@tempa{\textheight}{\dimexpr\ht\pandoc@box+\dp\pandoc@box\relax}%
  \Gscale@div\@tempb{\linewidth}{\wd\pandoc@box}%
  \ifdim\@tempb\p@<\@tempa\p@\let\@tempa\@tempb\fi% select the smaller of both
  \ifdim\@tempa\p@<\p@\scalebox{\@tempa}{\usebox\pandoc@box}%
  \else\usebox{\pandoc@box}%
  \fi%
}
% Set default figure placement to htbp
\def\fps@figure{htbp}
\makeatother
\setlength{\emergencystretch}{3em} % prevent overfull lines
\providecommand{\tightlist}{%
  \setlength{\itemsep}{0pt}\setlength{\parskip}{0pt}}
\usepackage{bookmark}
\IfFileExists{xurl.sty}{\usepackage{xurl}}{} % add URL line breaks if available
\urlstyle{same}
\hypersetup{
  pdftitle={Thanos benchmark report},
  pdfauthor={Thanos project (github.com/NathanSiemers/Thanos)},
  colorlinks=true,
  linkcolor={Maroon},
  filecolor={Maroon},
  citecolor={Blue},
  urlcolor={Blue},
  pdfcreator={LaTeX via pandoc}}

\title{Thanos benchmark report}
\usepackage{etoolbox}
\makeatletter
\providecommand{\subtitle}[1]{% add subtitle to \maketitle
  \apptocmd{\@title}{\par {\large #1 \par}}{}{}
}
\makeatother
\subtitle{Interactive cross-filtering for R/Shiny --- performance across
scales, backends, and storage layouts}
\author{Thanos project (github.com/NathanSiemers/Thanos)}
\date{2026-08-05}

\begin{document}
\maketitle

{
\setcounter{tocdepth}{2}
\tableofcontents
}
\section{What this report covers}\label{what-this-report-covers}

Thanos is a reusable R/Shiny module for interactive cross-filtering: the
user picks columns, gets auto-generated filter widgets, and every
selected column shows a live histogram of the rows passing all
\emph{other} filters, with that column's own selection overlaid. Every
slider drag or checkbox click must therefore recompute masks and repaint
plots --- performance is the product.

This report collects the project's benchmarks in the order they were
run, from a 337 thousand-row in-memory data frame to 38 million taxi
rows in DuckDB. Each section explains \textbf{what is being measured and
why it matters to interactivity}, then gives the numbers and their
interpretation. All results come from the scripts in \texttt{bench/} and
are reproducible:

\begin{verbatim}
Rscript bench/bench_masks.R      # the per-interaction data path
Rscript bench/bench_backends.R   # in-memory vs SQLite column store
Rscript bench/bench_big.R        # 38M rows: SQLite vs DuckDB aggregates
Rscript bench/bench_layouts.R    # melt vs wide vs parquet-direct
\end{verbatim}

Test machine: R 4.5.3, Linux, 28 cores, 377 GB RAM (rstudio container).
Times are means over repeated runs after a warm-up iteration; databases
are on local disk.

A note on architecture, since it frames every number below: the module
recomputes \textbf{only one column's mask} when that column's filter
changes, combines all masks with cumulative ANDs (O(3k) vector
operations for k selected columns, not O(k²)), and renders plots from
\textbf{pre-binned counts} (O(bins), not O(rows)). The benchmarks
measure exactly these steps.

\section{1. The per-interaction data path (flights, 336,776 rows, in
memory)}\label{the-per-interaction-data-path-flights-336776-rows-in-memory}

\textbf{What is measured.} One user interaction --- say, dragging the
\texttt{dep\_delay} slider --- triggers this pipeline: recompute that
column's logical mask, re-combine the leave-one-out masks for every
selected column, re-tabulate binned counts, and rebuild each histogram.
If the whole pipeline stays well under \textasciitilde150 ms per plot,
the UI feels instant. \texttt{bench/bench\_masks.R} times each stage in
isolation on \texttt{nycflights13::flights}.

\textbf{Why the ``OLD'' row exists.} The pre-merge prototypes rebuilt
each histogram by handing all \textasciitilde337k raw values to
\texttt{ggplot2::geom\_histogram} on every interaction. The final row
measures that legacy approach for contrast with the pre-binned design.

{\def\LTcaptype{none} % do not increment counter
\begin{longtable}[]{@{}ll@{}}
\toprule\noalign{}
operation & time \\
\midrule\noalign{}
\endhead
\bottomrule\noalign{}
\endlastfoot
make\_mask numeric (dep\_delay in {[}-10, 60{]}) & 10.7 ms \\
make\_mask categorical (3 carriers) & 22.2 ms \\
leave-one-out combine, k = 3 vars & 15.8 ms \\
leave-one-out combine, k = 6 vars & 44.4 ms \\
leave-one-out combine, k = 10 vars & 57.8 ms \\
tabulate counts, numeric 50 bins & 6.6 ms \\
plot\_histo build (ggplot object), numeric & 22.3 ms \\
plot\_histo build (ggplot object), categorical & 25.7 ms \\
plot\_histo build + render to png, numeric & 245 ms \\
OLD: geom\_histogram over 337k raw rows (build + render) & 920 ms \\
\end{longtable}
}

\textbf{Interpretation.}

\begin{itemize}
\tightlist
\item
  The \emph{data} side of an interaction is cheap: one mask recompute
  (\textasciitilde10--20 ms) plus one leave-one-out combine
  (\textasciitilde16--58 ms depending on how many columns are selected)
  plus \textasciitilde7 ms of counting per plot. Comfortably inside the
  interactivity budget.
\item
  Plot \textbf{rendering} dominates: \textasciitilde245 ms per
  histogram, of which only \textasciitilde22 ms is building the ggplot
  object --- the rest is fixed ggplot/grid/png-device overhead that does
  not grow with data size. Pre-binning made each render 3.7× faster than
  the legacy full-data approach (920 ms), and decoupled render cost from
  row count entirely.
\item
  Future lever if many panels ever feel sluggish: a faster graphics
  device (ragg) or a base-graphics \texttt{barplot} re-implementation of
  the same visual. Not currently needed.
\item
  Not visible in the table but architecturally decisive: the old module
  \emph{also} tore down and rebuilt every filter widget on every input
  (a \texttt{renderUI} dependency bug) and re-filtered the full data
  frame once per plot. The merged module touches one mask and k plots,
  nothing else.
\end{itemize}

\section{2. In-memory vs tall/skinny SQLite
(flights)}\label{in-memory-vs-tallskinny-sqlite-flights}

\textbf{What is measured.} The same module can run against an in-memory
data frame (\texttt{backend\_memory}) or a SQLite database in the
tall/skinny layout
\texttt{long\_data(row\_id,\ column\_name,\ value\_num,\ value\_txt)}
with a precomputed \texttt{column\_registry} (\texttt{backend\_sqlite}).
The module treats both as \textbf{column stores}: when the user selects
a column it is fetched once, in full, and all filtering then happens in
R on the cached vector. So the only structural difference between
backends is the cost of that one-time fetch, plus the cost of the
metadata needed to build widgets. \texttt{bench/bench\_backends.R}
measures both. Database: 6,352,149 long rows, 282 MB, indexed on
\texttt{(column\_name,\ row\_id)}.

{\def\LTcaptype{none} % do not increment counter
\begin{longtable}[]{@{}
  >{\raggedright\arraybackslash}p{(\linewidth - 4\tabcolsep) * \real{0.3333}}
  >{\raggedright\arraybackslash}p{(\linewidth - 4\tabcolsep) * \real{0.3333}}
  >{\raggedright\arraybackslash}p{(\linewidth - 4\tabcolsep) * \real{0.3333}}@{}}
\toprule\noalign{}
\begin{minipage}[b]{\linewidth}\raggedright
operation
\end{minipage} & \begin{minipage}[b]{\linewidth}\raggedright
memory
\end{minipage} & \begin{minipage}[b]{\linewidth}\raggedright
sqlite
\end{minipage} \\
\midrule\noalign{}
\endhead
\bottomrule\noalign{}
\endlastfoot
backend open (registry read) & --- & 340 ms \\
get\_column\_info, all 19 cols & 322 ms (first scan) & 2.2 ms
(registry) \\
get\_column(dep\_delay) & \textasciitilde0 ms & 203 ms \\
get\_column(distance) & \textasciitilde0 ms & 210 ms \\
get\_column(carrier) & \textasciitilde0 ms & 227 ms \\
get\_column(dest) & \textasciitilde0 ms & 231 ms \\
make\_mask on fetched column & 8.5 ms & 8.5 ms (identical) \\
\end{longtable}
}

\textbf{Interpretation.}

\begin{itemize}
\tightlist
\item
  The project hypothesis (``might not be too different, because this
  isn't a huge database'') is confirmed, and the benchmark shows
  \emph{why}: after a \textasciitilde200--230 ms one-time fetch per
  selected column, every per-interaction cost is byte-identical between
  backends, because the module filters cached vectors either way.
\item
  The registry pays for itself immediately: widget metadata (ranges,
  levels, NA counts, distinct-value counts, robust quantiles) is
  \textasciitilde146× faster from the registry than a first in-memory
  scan, because it was precomputed at build time.
\item
  The \texttt{(column\_name,\ row\_id)} index is what keeps fetches at
  \textasciitilde200 ms over a 6.4M-row long table; without it every
  fetch is a full scan.
\end{itemize}

\section{3. 38 million rows: SQLite vs DuckDB in aggregate
mode}\label{million-rows-sqlite-vs-duckdb-in-aggregate-mode}

\textbf{What is measured.} At taxi scale (NYC TLC yellow cab 2023:
38,310,226 rows, 13 columns kept, 460--494M cells in the melt) the
column-store approach stops being reasonable --- a single fetched column
is \textasciitilde300 MB, and R-side boolean algebra on 38M-long vectors
costs seconds per interaction. The module therefore switches to
\textbf{aggregate mode}: no column vector ever enters R; histogram
counts, row counts, and the row mask are SQL queries composed from the
current filter state, and only ≤ 50 bin counts cross the wire per plot.
\texttt{bench/bench\_big.R} runs those exact queries against the same
tall/skinny schema built in both engines, with two filters active
(\texttt{fare\_amount} in {[}5, 60{]}, \texttt{payment\_type} in \{1,
2\}) --- a realistic mid-interaction state.

\textbf{Build cost} (12 parquet months, \textasciitilde607 MB
compressed):

{\def\LTcaptype{none} % do not increment counter
\begin{longtable}[]{@{}lll@{}}
\toprule\noalign{}
step & SQLite (R melt via arrow/data.table) & DuckDB (all-SQL) \\
\midrule\noalign{}
\endhead
\bottomrule\noalign{}
\endlastfoot
build & 559 s & 94 s \\
index & 690 s & 64 s \\
file size & 27.8 GB (460M long rows*) & 5.9 GB (494M long rows) \\
\end{longtable}
}

* The SQLite build initially missed \texttt{airport\_fee} in months that
name the column \texttt{Airport\_fee}: R/arrow matches column names
case-sensitively where DuckDB does not. The build script now normalizes
the name; benchmark columns were unaffected.

\textbf{Per-interaction aggregate queries:}

{\def\LTcaptype{none} % do not increment counter
\begin{longtable}[]{@{}lll@{}}
\toprule\noalign{}
operation & DuckDB & SQLite \\
\midrule\noalign{}
\endhead
\bottomrule\noalign{}
\endlastfoot
get\_binned numeric (GROUP BY over filtered rows) & 477 ms & 75.5 s \\
get\_binned categorical & 229 ms & 50.7 s \\
get\_count & 361 ms & 45.1 s \\
get\_column full fetch (38M values) & 3.3 s & 28.3 s \\
get\_row\_mask & 0.9 s & 55.5 s \\
\end{longtable}
}

\textbf{Interpretation.}

\begin{itemize}
\tightlist
\item
  \textbf{DuckDB is the answer at this scale.} Sub-second aggregates
  make the module genuinely interactive on 38M rows: with the 500 ms
  debounce, a slider drag settles in roughly a second per histogram. Its
  columnar, parallel execution engine is well matched to the row-set
  semi-joins the tall/skinny layout requires.
\item
  \textbf{SQLite does not survive this scale in aggregate mode.} Every
  query walks hundreds of millions of b-tree index entries on a single
  core; 45--75 s per query is two orders of magnitude off. This is not
  an indictment of SQLite generally --- section 2 shows it is perfectly
  good as a column store up to roughly flights scale --- it is the wrong
  engine for large analytical scans.
\item
  The same module code drives both engines; the demo apps differ only in
  the backend constructor line. That was the point of the backend
  contract, and it is what made this an apples-to-apples comparison.
\end{itemize}

\section{4. Storage layouts: melt vs wide vs parquet-direct
(DuckDB)}\label{storage-layouts-melt-vs-wide-vs-parquet-direct-duckdb}

\textbf{What is measured.} Section 3 fixed the engine question; this
section asks whether the tall/skinny \textbf{layout} itself costs
performance. All three variants run in the \emph{same} DuckDB process
over the \emph{same} 38,310,226 rows, so the only variable is physical
layout:

\begin{itemize}
\tightlist
\item
  \textbf{melt} --- the production path:
  \texttt{long\_data(row\_id,\ column\_name,\ value\_num,\ value\_txt)}.
  Because there are no real columns, each filter becomes a
  \emph{semi-join}:
  \texttt{row\_id\ NOT\ IN\ (SELECT\ row\_id\ FROM\ long\_data\ WHERE\ column\_name\ =\ \textquotesingle{}fare\_amount\textquotesingle{}\ AND\ ...)}
  against the same 494M-row table. Column fetches ship
  \texttt{(row\_id,\ value)} pairs --- roughly twice the necessary bytes
  --- and are reassembled into full-length vectors in R (which is also
  how NA-by-absence is reconstructed).
\item
  \textbf{wide} --- one real column per variable (the \texttt{wide}
  table the DuckDB build materializes anyway as an intermediate).
  Filters are ordinary \texttt{WHERE\ fare\_amount\ BETWEEN\ 5\ AND\ 60}
  predicates: no joins. Fetches ship just the ordered column.
\item
  \textbf{parquet-direct} --- the same expressions evaluated straight
  over
  \texttt{read\_parquet(\textquotesingle{}yellow\_tripdata\_*.parquet\textquotesingle{})}
  with \textbf{no database build at all}. The one capability it lacks:
  parquet scan order is not guaranteed stable across queries, so there
  is no trustworthy row ID --- histogram counts and row counts are
  exact, but the \texttt{rows()}/\texttt{mask()} pointer a parent app
  consumes would require materializing row IDs first (at which point you
  have built the wide table).
\end{itemize}

Same interaction state as section 3.

{\def\LTcaptype{none} % do not increment counter
\begin{longtable}[]{@{}llll@{}}
\toprule\noalign{}
operation & melt & wide & parquet-direct \\
\midrule\noalign{}
\endhead
\bottomrule\noalign{}
\endlastfoot
get\_binned numeric, 2 filters & 477 ms & 168 ms & 633 ms \\
get\_binned categorical, 1 filter & 233 ms & 58 ms & 419 ms \\
get\_count, 2 filters & 358 ms & 76 ms & 402 ms \\
get\_column full fetch & 3.3 s & 1.6 s & 1.5 s \\
get\_row\_mask, 2 filters & 819 ms & 372 ms & n/a (no stable row\_id) \\
\end{longtable}
}

\textbf{Interpretation.}

\begin{itemize}
\tightlist
\item
  \textbf{The melt tax is real but bounded: roughly 3--5× on
  aggregates.} Wide's plain column predicates beat the melt's
  self-semi-joins on every query (58--168 ms vs 233--477 ms). DuckDB
  absorbs the join work impressively --- half a second over 494M long
  rows --- but it cannot make a join as cheap as no join.
\item
  \textbf{Parquet-direct is remarkable for what it doesn't need.} With
  zero build time, zero database file, and \textasciitilde600 MB of
  compressed source data, it lands in the same class as the melt (it
  pays decompression and file scanning per query instead of join work).
  For exploratory ``point Thanos at a directory of parquet and go'' use,
  it is viable today for histograms and counts.
\item
  \textbf{Practical guidance.} The tall/skinny melt remains the right
  \emph{interchange} schema --- it is what lets one
  \texttt{backend\_dbi} implementation serve SQLite and DuckDB
  identically, matches the project's relational-modeling goal, and its
  absolute numbers are interactive. When maximum speed matters, a
  wide-table DuckDB backend is a drop-in win (the backend contract hides
  the layout from the module entirely); the build even materializes the
  wide table already. A parquet-direct backend would additionally need
  row-ID materialization to support the parent-app row pointer.
\end{itemize}

\section{Overall guidance by scale}\label{overall-guidance-by-scale}

{\def\LTcaptype{none} % do not increment counter
\begin{longtable}[]{@{}
  >{\raggedright\arraybackslash}p{(\linewidth - 4\tabcolsep) * \real{0.3333}}
  >{\raggedright\arraybackslash}p{(\linewidth - 4\tabcolsep) * \real{0.3333}}
  >{\raggedright\arraybackslash}p{(\linewidth - 4\tabcolsep) * \real{0.3333}}@{}}
\toprule\noalign{}
\begin{minipage}[b]{\linewidth}\raggedright
data size
\end{minipage} & \begin{minipage}[b]{\linewidth}\raggedright
recommended setup
\end{minipage} & \begin{minipage}[b]{\linewidth}\raggedright
expected feel
\end{minipage} \\
\midrule\noalign{}
\endhead
\bottomrule\noalign{}
\endlastfoot
≤ \textasciitilde1M rows & \texttt{backend\_memory}, vector mode &
instant \\
\textasciitilde1M rows, data in a DB & \texttt{backend\_sqlite} column
store, vector mode & instant after \textasciitilde0.2 s per column
selection \\
tens of millions of rows & \texttt{backend\_duckdb} (melt), aggregate
mode & \textasciitilde0.5--1 s per interaction \\
same, maximum speed & wide-table DuckDB backend (future drop-in) &
\textasciitilde0.1--0.2 s per interaction \\
ad-hoc parquet exploration & parquet-direct backend (future; no row
pointer yet) & \textasciitilde0.5--1 s per interaction, zero build \\
\end{longtable}
}

\end{document}
